\documentclass[12pt,a4paper]{article}
\usepackage[margin=2.5cm]{geometry}
\usepackage{amsmath,amssymb,graphicx,hyperref,natbib,booktabs,caption,xcolor}
\definecolor{ai4sgreen}{HTML}{0D3B2E}
\definecolor{ai4sgold}{HTML}{C8956C}

\title{\textbf{\textcolor{ai4sgold}{AI4S-Journal} Method Article}}
\author{Author Name\textsuperscript{1*}, Co-Author\textsuperscript{2} \
\textsuperscript{1}Department, Institution, City, Country \
\textsuperscript{2}Department, Institution, City, Country \
\textsuperscript{*}Email: email@institution.edu}
\date{}

\begin{document}
\maketitle

\begin{abstract}
Describe the novel method, the problem it solves, and its key advantages
over existing approaches. Include quantitative performance highlights.
\end{abstract}
\noindent\textbf{Keywords:} algorithm, method, computational, [domain]

\section{Introduction}
Motivate the problem and review existing methods. State their limitations
and how your method addresses them.

\section{Method}
\subsection{Overview}
\subsection{Detailed Description}
Use formal notation where appropriate.

\subsection{Theoretical Analysis (optional)}
Convergence guarantees, complexity analysis, or generalization bounds.

\section{Experiments}
\subsection{Setup}
Describe datasets, baselines, hardware, and evaluation protocol.
\subsection{Results}
Quantitative comparisons with state-of-the-art methods.
\subsection{Ablation Studies}
Validate design choices through controlled experiments.

\section{Discussion}
\section{Conclusion}
\section*{Data and Code Availability}

\bibliographystyle{unsrt}
\begin{thebibliography}{99}
\bibitem{m1} Author, A. (Year). \textit{Title}. Conference/Journal.
\end{thebibliography}
\end{document}
